%% Table 3: Threat Landscape Summary (from §3, end)
%% Requires: \usepackage{tabularx}

\begin{table*}[htbp]
\centering
\caption{Threat landscape summary: type, affected modalities, severity, and defense maturity.}
\label{tab:threat-summary}
\scriptsize
\setlength{\tabcolsep}{4pt}
\renewcommand{\arraystretch}{0.92}
\begin{tabularx}{\textwidth}{l l c c X}
\toprule
\textbf{Threat Type} & \textbf{Modalities} & \textbf{Severity} & \textbf{Defense} & \textbf{Key Evidence} \\
\midrule
AI-generated phishing & Text, audio & Critical & Low & 82.6\% AI-generated; surpasses human red teams \\
Non-consensual synthetic & Image, video & Critical & Very low & 1,325\% increase in NCMEC reports; 440K in H1 2025 \\
Data poisoning / backdoors & All & High & Low & 250 documents suffice; $>$90\% ASR for instruction backdoors \\
Model extraction & All & High & Medium & Practical for \$60; industrial-scale distillation alleged \\
Prompt injection & Text & Critical & Very low & Structural problem; ``may never be fully solved'' \\
Jailbreaks & Text, image, audio & High & Low & $\sim$50\% success in 10 attempts; no uniform defense \\
Agentic exploitation & Multi-modal & Critical & Very low & 100\% of AI IDEs vulnerable; 72.8\% tool poisoning ASR \\
Deepfakes & Image, audio, video & Critical & Low & \$25M Arup fraud; 2,137\% surge in deepfake attacks \\
Information operations & Text, image, video & High & Low & AI disinfo perceived as more credible than human-written \\
\bottomrule
\end{tabularx}
\vspace{1pt}
\par\noindent{\tiny\emph{Note:} Defense maturity reflects deployment coverage, literature volume, and demonstrated effectiveness as of March 2026. Very low = emerging research, no deployed countermeasure; Low = active research, limited deployment; Medium = established methods with known limitations.}
\end{table*}
